% ===============================================================
%  جلسه بیست‌وششم
% ===============================================================
\chapter{آزمون‌های $t$، دونمونه‌ای، فیشر، $\chi^2$ و جایگشت}

این جلسه ادامه آزمون فرض است. اسلایدها ابتدا مراحل انجام آزمون فرض، تعاریف (فرضیه ساده و مرکب، آماره، فرضیه باطل و جایگزین)، خطاهای نوع اول و دوم و آزمون $z$ را مرور می‌کنند که به طور کامل در جلسه ۲۵ آمده است. مطالب جدید: \textbf{نمایش گرافیکی مقدار \lr{p}}، \textbf{آزمون‌های دونمونه‌ای}، \textbf{آزمون $t$} برای واریانس نامعلوم، \textbf{آزمون‌های یک‌طرفه و دوطرفه}، و سپس (به عنوان مطالب غیرامتحانی طبق اسلاید) \textbf{آزمون دقیق فیشر}، \textbf{آزمون $\chi^2$} و \textbf{آزمون جایگشت}. فصل با اشاره‌ای به \textbf{لم نیمن--پیرسون} تمام می‌شود.

(در فایل اسلاید این جلسه، عنوان «جلسه بیست‌وهفتم» درج شده است.)

\begin{summarybox}[خلاصه جلسه قبل]
\begin{itemize}
\item مراحل آزمون: فرض‌سازی ($H_0$ و $H_1$) $\leftarrow$ جمع‌آوری نمونه \lr{i.i.d.} $\leftarrow$ محاسبه آماره $\leftarrow$ مقایسه با ناحیه رد و گزارش احتمال.
\item $H_0$ (فرضیه باطل، معمولاً ساده): آنچه می‌خواهیم رد کنیم. $H_1$ (جایگزین، معمولاً مرکب): آنچه می‌خواهیم نشان دهیم.
\item خطای نوع اول (رد $H_0$ درست) با احتمال $\alpha$ = سطح معناداری؛ خطای نوع دوم (عدم رد $H_0$ نادرست) با احتمال $\beta$؛ توان $=1-\beta$.
\item مقدار \lr{p}: احتمال دیدن آماره‌ای حداقل به این اندازه پرت، اگر $H_0$ درست باشد. رد $H_0$ اگر و تنها اگر $\text{p}\le\alpha$.
\item آزمون $z$ ($\sigma$ معلوم): $z_n=\sqrt n(\bar X_n-\mu_0)/\sigma\sim\mathcal N(0,1)$ تحت $H_0$؛ رد اگر $|z_n|\ge z_{\alpha/2}$.
\end{itemize}
\end{summarybox}

\section{نمایش گرافیکی نحوه محاسبه مقدار \lr{p}}

\begin{center}
\begin{tikzpicture}
\end{tikzpicture}\\
{\small محور افقی: مجموعه نتایج ممکن (مقادیر آماره)؛ محور عمودی: احتمال (چگالی) مشاهده هر نتیجه \textbf{اگر $H_0$ درست باشد}. قله منحنی، مقدار واقعی تحت فرضیه باطل و محتمل‌ترین مشاهده است. خط‌چین، آستانه معناداری ۹۵٪ (یک‌طرفه) است. مقدار \lr{p} برای نتیجه مشاهده‌شده (نقطه سبز)، مساحت زیر منحنی از آن نقطه به بعد (ناحیه پررنگ) است. دو انتها، مشاهدات بسیار نامحتمل‌اند.}
\end{center}

شکل بالا تمام منطق آزمون فرض را در یک نگاه نشان می‌دهد:
\begin{itemize}
\item منحنی، توزیع آماره آزمون \textbf{با فرض درستی $H_0$} است؛ پس ابتدا باید بتوانیم این توزیع را (دقیق یا تقریبی) حساب کنیم. به همین دلیل آماره را طوری می‌سازیم که تحت $H_0$ توزیع شناخته‌شده‌ای داشته باشد ($\mathcal N(0,1)$، $t_{n-1}$، $\chi^2$، \ldots).
\item هرچه نتیجه مشاهده‌شده دورتر در دم باشد، مساحت دم پس از آن (مقدار \lr{p}) کمتر است و شاهد قوی‌تری علیه $H_0$ داریم.
\item اگر نقطه مشاهده‌شده از آستانه عبور کند، مساحت پس از آن کمتر از مساحت پس از آستانه ($0.05$) است؛ یعنی «عبور از آستانه» و «$\text{p}<0.05$» دو بیان یک چیزند.
\end{itemize}

\section{آزمون‌های تک‌نمونه‌ای و دونمونه‌ای \lr{(One-sample vs. Two-sample tests)}}
\begin{itemize}
\item در آزمون‌های \textbf{تک‌نمونه‌ای}، به دنبال (برای مثال) مقایسه میانگین یک توزیع با یک مقدار پیش‌فرض مانند $\mu_0$ هستیم که توسط فرضیه باطل \lr{(Null hypothesis)} مشخص می‌شود. (آزمون $z$ جلسه قبل از این نوع بود.)
\item در آزمون‌های \textbf{دونمونه‌ای}، از دو توزیع مختلف به مقادیر $n_1$ و $n_2$ نمونه تصادفی در دست بوده و به دنبال (برای مثال) آزمون این فرضیه هستیم که میانگین آنان با یکدیگر برابر است یا خیر.
\begin{itemize}
\item در این موارد، معمولاً از طریق تعریف یک آماره مناسب، آزمون دونمونه‌ای را به تک‌نمونه‌ای کاهش می‌دهیم.
\end{itemize}
\end{itemize}

\subsection{آزمون $z$ دونمونه‌ای}
فرض کنید $X_1,\ldots,X_{n_1}$ نمونه‌های \lr{i.i.d.} از توزیعی با میانگین $\mu_X$ و واریانس معلوم $\sigma_X^2$، و $Y_1,\ldots,Y_{n_2}$ نمونه‌های \lr{i.i.d.} (و مستقل از $X$ها) از توزیعی با میانگین $\mu_Y$ و واریانس معلوم $\sigma_Y^2$ باشند. فرضیه‌ها:
\[
H_0:\ \mu_X=\mu_Y\qquad\text{در برابر}\qquad H_1:\ \mu_X\neq\mu_Y .
\]
\textbf{کاهش به تک‌نمونه‌ای:} کمیت $D=\bar X_{n_1}-\bar Y_{n_2}$ را در نظر می‌گیریم. میانگین و واریانس آن (با استفاده از استقلال دو گروه):
\[
\E D=\mu_X-\mu_Y,\qquad \Var D=\Var\bar X_{n_1}+\Var\bar Y_{n_2}=\frac{\sigma_X^2}{n_1}+\frac{\sigma_Y^2}{n_2}.
\]
(واریانس تفاضل دو متغیر \textbf{مستقل}، \textbf{مجموع} واریانس‌هاست، نه تفاضل آن‌ها: هر دو منبع نوسان به عدم قطعیت $D$ اضافه می‌کنند.) تحت $H_0$ داریم $\E D=0$؛ پس آماره
\[
\boxed{\ z=\frac{\bar X_{n_1}-\bar Y_{n_2}}{\sqrt{\dfrac{\sigma_X^2}{n_1}+\dfrac{\sigma_Y^2}{n_2}}}\ }
\]
تحت $H_0$ (دقیقاً برای داده گاوسی، و تقریباً طبق حد مرکزی) توزیع $\mathcal N(0,1)$ دارد و دقیقاً مثل آزمون $z$ تک‌نمونه‌ای عمل می‌کنیم.

\begin{example}[آیا زنان و مردان در آزمون هوش متفاوت‌اند؟]
اسلاید این پرسش را مثال می‌زند: «آیا زنان در آزمون \lr{IQ} نمره‌ای به طور معنادار متفاوت از مردان می‌گیرند؟» فرض کنید نمره \lr{IQ} در هر دو گروه انحراف معیار معلوم $\sigma=15$ دارد. از هر گروه $50$ نفر آزمون داده‌اند: میانگین زنان $102.5$ و میانگین مردان $98.0$. آزمون دوطرفه را انجام دهید.
\begin{solution}
\[
\sqrt{\frac{15^2}{50}+\frac{15^2}{50}}=\sqrt{4.5+4.5}=3,\qquad z=\frac{102.5-98.0}{3}=1.5 .
\]
مقدار \lr{p}: $2\big(1-\Phi(1.5)\big)=2\times0.0668=0.134$. در سطح ۵٪، $H_0$ رد نمی‌شود: اختلاف $4.5$ واحدی در نمونه‌هایی با این اندازه، به‌راحتی می‌تواند حاصل تصادف باشد. (توجه: «رد نشدن» به معنی اثبات برابری نیست؛ شاید با نمونه بزرگ‌تر اختلاف معنادار شود.)
\end{solution}
\end{example}

\begin{note}
\notetitle \textbf{نمونه‌های جفتی \lr{(paired)}.} اگر دو نمونه مستقل نباشند بلکه \textbf{جفت‌جفت} باشند (مثلاً فشار خون هر بیمار قبل و بعد از دارو)، بهتر است تفاضل هر جفت $D_i=X_i-Y_i$ را حساب کنیم و آزمون \textbf{تک‌نمونه‌ای} $H_0:\ \E D=0$ را روی $D_i$ها انجام دهیم. این روش تفاوت‌های فردی (که در هر دو اندازه‌گیری مشترک‌اند) را حذف می‌کند و معمولاً توان بسیار بیشتری دارد. (تمرین ۱ را ببینید.)
\end{note}

\section{آزمون میانگین نمونه‌ای با واریانس نامعلوم \lr{(t-test)}}

\subsection{صورت مسئله و آماره}
\begin{itemize}
\item فرض کنید که نمونه‌های $X_1,X_2,\ldots,X_n$ به صورت \lr{i.i.d.} از یک توزیع نامعلوم تولید شده باشند. فرض کنید که انحراف معیار توزیع دانسته شده \textbf{نیست}.
\item هدف طراحی آزمون این فرضیه است که میانگین توزیع مقدار $\mu_0$ را دارا است یا خیر:
\[
H_0:\ \E(X)=\mu_0,\qquad H_1:\ \E(X)\neq\mu_0 .
\]
\end{itemize}
آماره $z_n$ قابل محاسبه نیست، چون $\sigma$ را نداریم. مانند جلسه ۲۳، $\sigma$ را با تخمین آن جایگزین می‌کنیم و آماره زیر را به عنوان آماره آزمون انتخاب می‌کنیم:
\[
\boxed{\ t_n=\frac{\sqrt n\,\big(\bar X_n-\mu_0\big)}{\bar\sigma_n}\ },\qquad
\bar\sigma_n^2=\frac{1}{n-1}\sum_{i=1}^n\big(X_i-\bar X_n\big)^2,
\]
که در اینجا مقصود از $\bar\sigma_n$ \textbf{انحراف معیار نمونه‌ای} است (همان $\bar S_n$ جلسه ۲۳؛ در اسلاید در فرمول واریانس، خط بالای $\sigma_n^2$ جا افتاده است).

\subsection{توزیع آماره}
\begin{itemize}
\item اگر $X_i$ها که نمونه‌های محقق‌شده $X$ هستند توزیع گاوسی داشته باشند، مقدار $t_n$ از یک توزیع \lr{Student's t-distribution} با $n-1$ درجه آزادی پیروی نموده که ممکن است شیفت یافته باشد، و مقدار شیفت به برقراری فرضیه باطل $H_0$ یا فرضیه جایگزین بستگی دارد.
\end{itemize}
دقیق‌تر: \textbf{تحت $H_0$} (یعنی وقتی $\mu_0$ همان میانگین واقعی است)، $t_n$ دقیقاً همان کمیت محوری $\bar T_n$ جلسه ۲۳ است و
\[
t_n\ \sim\ t_{n-1}\qquad(\text{تحت}\ H_0).
\]
تحت $H_1$، صورت کسر شامل جمله اضافه $\sqrt n(\mu-\mu_0)$ است که آماره را از صفر دور می‌کند (توزیع حاصل را «$t$ غیرمرکزی» می‌نامند و با افزایش $n$ مانند $\sqrt n$ از صفر دور می‌شود).
\begin{itemize}
\item اگر $X$ توزیع گاوسی نداشته باشد نیز، با افزایش محدود $n$، توزیع $t_n$ به یک توزیع \lr{t-distribution} نزدیک و نزدیک‌تر می‌گردد. (دلیل: طبق حد مرکزی صورت تقریباً گاوسی و طبق قانون اعداد بزرگ $\bar\sigma_n\to\sigma$؛ پس $t_n$ تقریباً $\mathcal N(0,1)$ است و $t_{n-1}$ هم برای $n$ بزرگ تقریباً همان $\mathcal N(0,1)$ است.)
\item سپس، می‌توان مانند مورد $z$-\lr{test} عمل کرده و مقدار \lr{p-value} را از روی دم توزیع محاسبه نمود. تنها تفاوت، استفاده از \lr{CDF} توزیع \lr{t-distribution} به جای توزیع نرمال است:
\[
\text{p-value}=\Prob\big(|T|\ge|t_{\text{obs}}|\big)=2\Big(1-F_{t,n-1}\big(|t_{\text{obs}}|\big)\Big),\qquad\text{رد}\ H_0\ \text{اگر}\ |t_{\text{obs}}|\ge t_{n-1,\alpha/2}.
\]
\end{itemize}

\begin{example}[آزمون $t$: طول قطعه]
در مثال جلسه ۲۳ (ده اندازه‌گیری با $\bar x=10.05$ و $\bar s=0.3375$)، طول اسمی قطعه طبق نقشه $\mu_0=10.3$ میلی‌متر است. آیا شواهدی وجود دارد که میانگین طول قطعه‌ها با طول اسمی متفاوت است؟ ($\alpha=0.05$)
\begin{solution}
$H_0:\ \mu=10.3$ در برابر $H_1:\ \mu\neq10.3$. آماره:
\[
t_{\text{obs}}=\frac{\sqrt{10}\,(10.05-10.3)}{0.3375}=\frac{3.162\times(-0.25)}{0.3375}=-2.343 .
\]
درجه آزادی $9$ و مقدار بحرانی $t_{9,0.025}=2.262$. چون $|t_{\text{obs}}|=2.343>2.262$، $H_0$ رد می‌شود. مقدار \lr{p} (با نرم‌افزار یا جدول دقیق‌تر) حدود $0.044$ است. سازگاری با جلسه ۲۳: بازه اطمینان ۹۵٪ برابر $[9.81,\ 10.29]$ بود که $10.3$ را شامل نمی‌شود. اگر به اشتباه از آزمون $z$ با $\sigma=0.3375$ استفاده می‌کردیم، مقدار \lr{p} برابر $0.019$ می‌شد که شواهد را قوی‌تر از واقع نشان می‌دهد.
\end{solution}
\end{example}

\begin{note}
\notetitle \textbf{آزمون $t$ دونمونه‌ای.} برای مقایسه میانگین دو گروه مستقل با واریانس نامعلوم \textbf{ولی برابر}، واریانس مشترک را از ترکیب دو نمونه تخمین می‌زنیم:
\[
S_p^2=\frac{(n_1-1)S_X^2+(n_2-1)S_Y^2}{n_1+n_2-2},\qquad
t=\frac{\bar X-\bar Y}{S_p\sqrt{\frac1{n_1}+\frac1{n_2}}}\ \sim\ t_{n_1+n_2-2}\ \ (\text{تحت}\ H_0).
\]
(درجه آزادی: $n_1+n_2$ داده منهای ۲ میانگین تخمین‌زده‌شده.) اگر واریانس‌ها برابر نباشند، از نسخه تقریبی ولش \lr{(Welch)} استفاده می‌شود.
\end{note}

\section{آزمون‌های یک‌طرفه و دوطرفه \lr{(One-tailed vs. Two-tailed tests)}}
\begin{itemize}
\item گاهی اوقات توزیع آماره دو دم \lr{(tail)} داشته و هر دو کران متناظر با دوری \lr{(deviation)} از فرض باطل هستند، اما گاهی اوقات توزیع آماره یک کران دارد.
\end{itemize}

\begin{center}
\begin{tikzpicture}
\end{tikzpicture}\\
{\small راست: توزیع $t$ (متقارن)؛ هر دو دم ناحیه رد هستند (آزمون دوطرفه). چپ: توزیع $F$ (نامتقارن و نامنفی)؛ فقط دم راست آن ناحیه رد است (آزمون یک‌طرفه).}
\end{center}

دو موقعیت متفاوت به آزمون یک‌طرفه منجر می‌شود:
\begin{enumerate}
\item \textbf{ماهیت آماره:} بعضی آماره‌ها ذاتاً نامنفی‌اند و انحراف از $H_0$ فقط آن‌ها را \textbf{بزرگ} می‌کند. مثلاً آماره $\chi^2$ (در ادامه) مجموع مربعات انحراف‌هاست؛ هر نوع ناسازگاری با $H_0$ آن را بزرگ می‌کند و مقادیر نزدیک صفر یعنی سازگاری کامل. پس فقط دم راست ناحیه رد است. آماره $F$ (نسبت دو واریانس نمونه‌ای، برای مقایسه واریانس‌ها یا در تحلیل واریانس) هم چنین است.
\item \textbf{جهت‌دار بودن $H_1$:} حتی با آماره متقارن مثل $z$ یا $t$، اگر $H_1$ جهت‌دار باشد (مثلاً $\mu<\mu_0$ در مثال برنامه تلویزیونی)، فقط یک دم شاهد علیه $H_0$ است و کل $\alpha$ در همان دم قرار می‌گیرد: آستانه $z_\alpha$ (مثلاً $1.645$) به جای $z_{\alpha/2}$ (مثلاً $1.96$).
\end{enumerate}

\begin{warning}
\warntitle انتخاب یک‌طرفه یا دوطرفه بودن آزمون باید \textbf{پیش از دیدن داده‌ها} و بر اساس پرسش پژوهش انجام شود. اگر پس از دیدن جهت اثر، آزمون یک‌طرفه را در همان جهت انتخاب کنیم، در واقع مقدار \lr{p} را به ناحق نصف کرده‌ایم و احتمال خطای نوع اول واقعی دو برابر $\alpha$ اعلام‌شده می‌شود.
\end{warning}

\section{آزمون دقیق فیشر \lr{(Fisher's Exact Test)}}
\textbf{از اینجا به بعد (طبق اسلاید) غیرامتحانی است.}

\begin{itemize}
\item در بررسی‌های آماری، در موارد بسیاری به دنبال این هستیم که آیا ارتباطی میان دو ویژگی \textbf{دودویی} وجود دارد یا خیر؟
\begin{itemize}
\item برای مثال: آیا یک طبقه‌بندی‌کننده دودویی، به طرز معنی‌داری دو کلاس خاص را طبقه‌بندی می‌کند یا خیر؟ یا:
\item نمونه‌ای از نوجوانان را از یک طرف به دختر و پسر، و از طرف دیگر به کسانی که در حال درس خواندن برای امتحان آمار هستند و کسانی که نیستند تقسیم می‌کنیم. فرض ما این است که نسبت درس‌خوان‌ها در میان دختران بیشتر از پسران است، و می‌خواهیم آزمون کنیم که آیا اختلاف نسبت‌های مشاهده‌شده معنادار است. داده‌ها ممکن است چنین باشند:
\end{itemize}
\end{itemize}

\begin{center}
\renewcommand{\arraystretch}{1.4}
\begin{tabular}{c|c|c|c}
& پسران & دختران & جمع سطر\\\hline
در حال درس خواندن & $1$ & $9$ & $10$\\
درس نمی‌خوانند & $11$ & $3$ & $14$\\\hline
جمع ستون & $12$ & $12$ & $24$
\end{tabular}
\end{center}

\begin{itemize}
\item آزمون فیشر با جداولی موسوم به \textbf{جدول توافقی \lr{(Contingency Table)}} سروکار دارد. هدف این است که مشخص شود دو متغیر تصادفی دودویی که متناظر سطرها و ستون‌های این جدول هستند، از یکدیگر استقلال آماری دارند یا خیر.
\item این آزمون \textbf{دقیق \lr{(Exact)}} است؛ بدین معنا که مانند $z$-\lr{test} و یا $t$-\lr{test} از فرض گاوسی بودن (که در $n$های بزرگ به تدریج محقق می‌گردد) استفاده نشده و مقدار \lr{p-value} محاسبه‌شده به ازای هر اندازه از جامعه نمونه دقیق خواهد بود.
\end{itemize}

\subsection{منطق و محاسبه (روابط کامل)}
اسلاید برای روابط به صفحه ویکی‌پدیای این آزمون ارجاع می‌دهد؛ در اینجا آن‌ها را گام به گام می‌سازیم.

\textbf{گام ۱: شرطی کردن روی جمع‌های حاشیه‌ای.} جمع سطرها و ستون‌ها (۱۰ درس‌خوان، ۱۴ غیردرس‌خوان، ۱۲ پسر، ۱۲ دختر) را ثابت فرض می‌کنیم. با این جمع‌ها، کل جدول با \textbf{یک} عدد تعیین می‌شود: تعداد پسران درس‌خوان $a$ (بقیه خانه‌ها از تفاضل به دست می‌آیند).

\textbf{گام ۲: توزیع $a$ تحت $H_0$.} اگر جنسیت و درس خواندن مستقل باشند، اینکه کدام ۱۰ نفر از ۲۴ نفر درس‌خوان باشند هیچ ربطی به جنسیت ندارد. مثل این است که از کیسه‌ای شامل ۱۲ پسر و ۱۲ دختر، ۱۰ نفر را به تصادف (بدون جایگذاری) بیرون بکشیم و «درس‌خوان» بنامیم. تعداد پسران در این ۱۰ نفر توزیع \textbf{فوق‌هندسی} دارد: تعداد کل حالت‌های انتخاب ۱۰ نفر از ۲۴ نفر $\binom{24}{10}$ است و تعداد حالت‌هایی که دقیقاً $a$ نفر از آنان پسر (و $10-a$ نفر دختر) باشند $\binom{12}{a}\binom{12}{10-a}$؛ و چون همه انتخاب‌ها هم‌احتمال‌اند:
\[
\Prob(a)=\frac{\binom{12}{a}\binom{12}{10-a}}{\binom{24}{10}}
=\frac{\binom{a+b}{a}\binom{c+d}{c}}{\binom{n}{a+c}},
\]
که در فرم دوم (فرم استاندارد ویکی‌پدیا) $a,b$ سطر اول، $c,d$ سطر دوم و $n$ کل است (هر دو فرم یک مقدار می‌دهند).

\textbf{گام ۳: مقدار \lr{p}.} با $\binom{24}{12}=2{,}704{,}156$ و فرم دوم:
\[
\Prob(a=1)=\frac{\binom{10}{1}\binom{14}{11}}{\binom{24}{12}}=\frac{10\times364}{2704156}=0.001346,\qquad
\Prob(a=0)=\frac{\binom{10}{0}\binom{14}{12}}{\binom{24}{12}}=\frac{91}{2704156}=0.0000337 .
\]
\begin{itemize}
\item \textbf{یک‌طرفه} ($H_1$: درس‌خوان‌ها در دختران بیشترند، یعنی $a$ کوچک): جداولی حداقل به این اندازه افراطی، $a\le1$ هستند: $\text{p}=0.001346+0.0000337=0.00138$.
\item \textbf{دوطرفه:} به دلیل تقارن (۱۲ پسر و ۱۲ دختر)، جداول $a=9$ و $a=10$ هم به همان اندازه افراطی‌اند: $\text{p}=2\times0.00138=0.0028$.
\end{itemize}
پس ارتباط جنسیت و درس خواندن در این داده‌ها بسیار معنادار است.

\section{سایر آزمون‌ها}

\subsection{آزمون بهنجارش $\chi^2$ \lr{(Chi-squared test)}}
\begin{itemize}
\item فرض کنید که $n$ شیء به طریقی تصادفی و مستقل از یکدیگر به $k$ دسته \lr{(category)} تقسیم شده‌اند. فرضیه باطل (صفر) این است که احتمال تعلق به دسته $i$ام، برابر با $p_i$ است.
\item هدف سنجش صحت این فرضیه است.
\end{itemize}
فرض کنید $n$ مشاهده از یک نمونه تصادفی به $k$ دسته ناسازگار با تعدادهای مشاهده‌شده $x_i$ (برای $i=1,\ldots,k$) تقسیم شده‌اند و فرضیه باطل احتمال $p_i$ را برای افتادن هر مشاهده در دسته $i$ام می‌دهد. پس تعدادهای \textbf{مورد انتظار} برای همه $i$ها $m_i=np_i$ هستند، که
\[
\sum_{i=1}^kp_i=1,\qquad \sum_{i=1}^km_i=n\sum_{i=1}^kp_i=\sum_{i=1}^kx_i=n .
\]
پیرسون نشان داد که تحت درستی فرضیه باطل، وقتی $n\to\infty$، توزیع حدی کمیت زیر توزیع $\chi^2$ است:
\[
\boxed{\ X^2=\sum_{i=1}^k\frac{(x_i-m_i)^2}{m_i}=\sum_{i=1}^k\frac{x_i^2}{m_i}-n\ }\qquad\xrightarrow{\ n\to\infty\ }\qquad\chi^2_{k-1}.
\]
\textbf{اثبات تساوی دوم:} با بسط مربع
\[
\sum_i\frac{(x_i-m_i)^2}{m_i}=\sum_i\frac{x_i^2}{m_i}-2\sum_ix_i+\sum_im_i=\sum_i\frac{x_i^2}{m_i}-2n+n=\sum_i\frac{x_i^2}{m_i}-n .
\]
\textbf{شهود توزیع:} $x_i$ تعداد موفقیت در $n$ آزمایش با احتمال $p_i$ است (دوجمله‌ای)، پس تقریباً گاوسی با میانگین $m_i$ و واریانس $np_i(1-p_i)\approx m_i$. بنابراین $(x_i-m_i)/\sqrt{m_i}$ تقریباً گاوسی استاندارد است و $X^2$ مجموع $k$ مربع از چنین کمیت‌هایی است، یعنی تقریباً $\chi^2$. اما این $k$ کمیت آزاد نیستند: همواره $\sum_i(x_i-m_i)=0$. این یک قید خطی یک درجه آزادی را کم می‌کند و $k-1$ می‌ماند (همان استدلال $n-1$ در واریانس نمونه‌ای).

\textbf{تصمیم:} $X^2$ بزرگ یعنی تعدادهای مشاهده‌شده از تعدادهای مورد انتظار دورند؛ پس آزمون \textbf{یک‌طرفه} (دم راست) است: $H_0$ را رد می‌کنیم اگر $X^2$ از چندک $1-\alpha$ توزیع $\chi^2_{k-1}$ بیشتر باشد. قاعده سرانگشتی برای معتبر بودن تقریب: همه $m_i\ge5$.

\begin{example}[آیا تاس سالم است؟]
تاسی $120$ بار ریخته شده و تعداد آمدن وجه‌های ۱ تا ۶ به ترتیب $25,17,15,23,24,16$ بوده است. آیا تاس سالم است؟ ($\alpha=0.05$)
\begin{solution}
$H_0$: $p_i=\frac16$ برای همه وجه‌ها؛ پس $m_i=120/6=20$ (همگی $\ge5$). انحراف‌ها: $5,-3,-5,3,4,-4$ (مجموعشان صفر)، مربع‌ها: $25,9,25,9,16,16$ با مجموع $100$:
\[
X^2=\frac{100}{20}=5.0 .
\]
درجه آزادی $k-1=5$ و مقدار بحرانی $\chi^2_{5,0.95}=11.07$. چون $5.0<11.07$، $H_0$ رد نمی‌شود (مقدار \lr{p} حدود $0.42$). اختلاف‌ها در حد نوسانات تصادفی عادی‌اند. (برای مقایسه: میانگین $\chi^2_5$ برابر $5$ است؛ آماره مشاهده‌شده دقیقاً «معمولی» است.)
\end{solution}
\end{example}

\begin{note}
\notetitle \textbf{آزمون $\chi^2$ استقلال.} برای جدول توافقی با $r$ سطر و $c$ ستون، تعداد مورد انتظار هر خانه تحت استقلال $m_{ij}=\frac{(\text{جمع سطر}\ i)(\text{جمع ستون}\ j)}{n}$ است و $X^2=\sum_{i,j}\frac{(x_{ij}-m_{ij})^2}{m_{ij}}$ تقریباً $\chi^2_{(r-1)(c-1)}$ دارد. برای جدول نوجوانان بالا: $m=\begin{bmatrix}5&5\\7&7\end{bmatrix}$ و
\[
X^2=\frac{(1-5)^2}{5}+\frac{(9-5)^2}{5}+\frac{(11-7)^2}{7}+\frac{(3-7)^2}{7}=3.2+3.2+2.29+2.29=10.97,
\]
با یک درجه آزادی و مقدار \lr{p} تقریبی $0.0009$. این عدد حدود سه برابر کوچک‌تر از مقدار \lr{p} دقیق فیشر ($0.0028$) است: با تعدادهای کوچک، تقریب $\chi^2$ معناداری را بیش از واقع نشان می‌دهد؛ دقیقاً به همین دلیل آزمون دقیق فیشر ارزشمند است.
\end{note}

\subsection{آزمون جایگشت \lr{(Permutation Test)}}
\begin{itemize}
\item این آزمون کاربردهای فراوانی در عمل دارد.
\item فرض کنید که از دو دسته مشخص به شما مقداری نمونه تصادفی داده شده و قصد دارید بدانید که آیا میان این دو دسته داده، تفاوت معنی‌داری وجود دارد یا خیر.
\item مراحل:
\begin{enumerate}
\item ابتدا یک آماره را مشخص کنید (مثلاً تفاضل میانگین دو دسته). سپس تمام داده‌ها را روی هم بریزید و فعلاً دسته‌بندی اولیه داده‌شده را فراموش کنید.
\item سپس، به طرق تصادفی و به «تعداد کافی» دسته‌بندی تصادفی میان دادگان تجمیع‌شده ایجاد کرده و مقدار آماره را محاسبه کنید. سعی کنید که توزیع تجربی آماره را تحت این جایگشت‌ها پیدا کنید.
\item در آخر، مقدار آماره را با دسته‌بندی اولیه داده‌شده محاسبه کرده و مقدار \lr{p-value} را تخمین بزنید: کسری از جایگشت‌ها که آماره‌شان حداقل به اندازه آماره واقعی افراطی است.
\end{enumerate}
\end{itemize}
\textbf{چرا درست کار می‌کند؟} اگر $H_0$ درست باشد و هر دو دسته از یک توزیع آمده باشند، برچسب دسته‌ها هیچ اطلاعاتی ندارد: هر نحوه نسبت دادن برچسب‌ها به داده‌ها به یک اندازه محتمل است (\textbf{تعویض‌پذیری}). پس آماره واقعی صرفاً یک نمونه تصادفی از «توزیع جایگشتی» است و می‌توانیم ببینیم چقدر در دم آن قرار دارد. نکته زیبا: هیچ فرضی درباره شکل توزیع (گاوسی و غیره) لازم نیست.

\begin{example}[آزمون جایگشت دستی]
دسته الف: $12,15,18$ و دسته ب: $9,10,11$. آماره: تفاضل میانگین‌ها. آزمون جایگشت دقیق (همه جایگشت‌ها) را انجام دهید و با آزمون $t$ دونمونه‌ای مقایسه کنید.
\begin{solution}
آماره واقعی: $15-10=5$. شش داده تجمیع‌شده $9,10,11,12,15,18$ با مجموع $75$ را به $\binom63=20$ طریق می‌توان به دو دسته سه‌تایی تقسیم کرد. اگر مجموع دسته الف $S$ باشد، تفاضل میانگین‌ها $\frac{S-(75-S)}{3}=\frac{2S-75}{3}$ است. تفاضل $\ge5$ یعنی $S\ge45$؛ و تنها دسته‌ای با مجموع $45$، دقیقاً $\{12,15,18\}$ (سه عدد بزرگ) است. پس:
\[
\text{p}_{\text{یک‌طرفه}}=\frac{1}{20}=0.05,\qquad\text{p}_{\text{دوطرفه}}=\frac{2}{20}=0.10
\]
(در حالت دوطرفه، دسته $\{9,10,11\}$ با تفاضل $-5$ هم به همان اندازه افراطی است). نکته: با ۳+۳ داده، کوچک‌ترین مقدار \lr{p} دوطرفه ممکن $0.1$ است؛ آزمون جایگشت هرگز ادعایی قوی‌تر از آنچه داده‌ها اجازه می‌دهند نمی‌کند.

\textbf{مقایسه با $t$:} $S_X^2=9$، $S_Y^2=1$، $S_p^2=\frac{2\times9+2\times1}{4}=5$، $t=\frac{5}{\sqrt5\sqrt{2/3}}=2.74$ با ۴ درجه آزادی و $\text{p}\approx0.052$. آزمون $t$ (با فرض گاوسی بودن و واریانس برابر، که با واریانس‌های $9$ و $1$ مشکوک است) شواهد را قوی‌تر نشان می‌دهد.
\end{solution}
\end{example}

\begin{note}
\notetitle وقتی تعداد جایگشت‌ها بسیار زیاد است (مثلاً $\binom{40}{20}\approx1.4\times10^{11}$)، به جای همه آن‌ها چند هزار جایگشت تصادفی تولید می‌کنیم (روش مونت‌کارلو). کسر جایگشت‌های افراطی، تخمینی نااریب از مقدار \lr{p} است و دقت آن طبق قانون اعداد بزرگ با تعداد جایگشت‌ها بهتر می‌شود.
\end{note}

\section{مطالب جانبی و تکمیلی: لم نیمن--پیرسون}
اسلاید پیشنهاد می‌کند: «لم \lr{Neyman--Pearson} در رابطه با نرخ افت خطا در آزمون‌های آماری با افزایش نمونه را مطالعه کنید.» خلاصه آن را اینجا می‌آوریم.

\begin{theorem}[لم نیمن--پیرسون]
برای دو فرضیه \textbf{ساده} $H_0:\ X_i\sim p$ و $H_1:\ X_i\sim q$ (نمونه‌های \lr{i.i.d.})، در میان همه آزمون‌هایی که احتمال خطای نوع اول آن‌ها حداکثر $\alpha$ است، آزمونی که \textbf{بیشترین توان} را دارد، \textbf{آزمون نسبت درست‌نمایی} است:
\[
\text{رد}\ H_0\quad\Longleftrightarrow\quad\Lambda_n=\prod_{i=1}^n\frac{q(X_i)}{p(X_i)}>\eta,
\]
که آستانه $\eta$ طوری انتخاب می‌شود که $\Prob(\Lambda_n>\eta\given H_0)=\alpha$.
\end{theorem}
\textbf{شهود:} هر نقطه داده $x$ را باید به ناحیه رد اضافه کنیم اگر «ارزش» آن برای توان ($q(x)$) نسبت به «هزینه» آن در خطای نوع اول ($p(x)$) زیاد باشد؛ مثل پر کردن کوله‌پشتی با اقلامی که بیشترین نسبت ارزش به وزن را دارند. (این همان پرسش «$X\sim P$ یا $X\sim Q$؟» جلسه ۲۵ است.)

\textbf{نمونه:} برای $p=\mathcal N(0,1)$ و $q=\mathcal N(1,1)$:
\[
\log\Lambda_n=\sum_{i=1}^n\Big[-\frac{(X_i-1)^2}{2}+\frac{X_i^2}{2}\Big]=\sum_{i=1}^n\Big(X_i-\frac12\Big)=n\bar X_n-\frac n2,
\]
پس آزمون بهینه «$\bar X_n>c$» است: همان آزمون $z$ یک‌طرفه! یعنی آزمون $z$ برای این مسئله بهترین آزمون ممکن است.

\textbf{نرخ افت خطا:} با ثابت نگه داشتن $\alpha$، احتمال خطای نوع دوم بهترین آزمون با افزایش $n$ به صورت \textbf{نمایی} کاهش می‌یابد:
\[
\beta_n\approx e^{-n\,D(p\|q)},\qquad D(p\|q)=\E_p\Big[\log\frac{p(X)}{q(X)}\Big],
\]
که $D(p\|q)$ \textbf{واگرایی کولبک--لایبلر} (فاصله آماری میان دو توزیع) است (لم چرنوف--استاین). شهود: طبق قانون اعداد بزرگ، $\frac1n\log\Lambda_n$ تحت $H_0$ به $-D(p\|q)$ و تحت $H_1$ به $D(q\|p)$ همگرا می‌شود؛ هرچه دو توزیع از هم دورتر باشند، تمایزشان آسان‌تر و نرخ افت خطا بیشتر است. برای مثال گاوسی بالا $D(p\|q)=\frac12$.

\section{تمرین‌های تکمیلی}

\begin{exercise}[آزمون $t$ جفتی]
فشار خون ۶ بیمار قبل و بعد از مصرف یک دارو اندازه‌گیری شده و تفاضل‌ها (بعد منهای قبل) چنین است: $-8,-3,-6,2,-5,-4$. آیا دارو اثر دارد؟ آزمون دوطرفه با $\alpha=0.05$ انجام دهید.
\end{exercise}
\begin{solution}
روی تفاضل‌ها $D_i$ آزمون تک‌نمونه‌ای $H_0:\ \E D=0$ را انجام می‌دهیم. $\bar d=\frac{-24}{6}=-4$. انحراف‌ها از میانگین: $-4,1,-2,6,-1,0$ با مجموع مربعات $16+1+4+36+1+0=58$؛ پس $\bar s^2=58/5=11.6$ و $\bar s=3.406$.
\[
t=\frac{\sqrt6\times(-4)}{3.406}=\frac{2.449\times(-4)}{3.406}=-2.877 .
\]
با ۵ درجه آزادی، $t_{5,0.025}=2.571$ و $|t|>2.571$؛ پس $H_0$ رد می‌شود ($\text{p}\approx0.035$). اگر به جای آزمون جفتی، اندازه‌های قبل و بعد را دو گروه مستقل فرض می‌کردیم، تفاوت‌های زیاد میان بیماران (که در تفاضل حذف می‌شوند) واریانس را بزرگ می‌کرد و احتمالاً اثر آشکار نمی‌شد.
\end{solution}

\begin{exercise}[داده‌های مندل]
مندل در آزمایش نخودفرنگی، $705$ گل ارغوانی و $224$ گل سفید مشاهده کرد. نظریه او نسبت $3:1$ را پیش‌بینی می‌کند. با آزمون $\chi^2$ این نظریه را بسنجید.
\end{exercise}
\begin{solution}
$n=929$، $p_1=\frac34$، $p_2=\frac14$؛ پس $m_1=696.75$ و $m_2=232.25$. انحراف‌ها $\pm8.25$:
\[
X^2=\frac{8.25^2}{696.75}+\frac{8.25^2}{232.25}=0.098+0.293=0.391 .
\]
با $k-1=1$ درجه آزادی (مقدار بحرانی $3.84$)، $H_0$ رد نمی‌شود ($\text{p}\approx0.53$): داده‌ها با نظریه سازگارند.
\end{solution}

\begin{exercise}[طبقه‌بند با ۸۰٪ دقت]
یک طبقه‌بند روی ۱۰ نمونه (۵ از هر کلاس) آزمایش شده: از ۵ نمونه کلاس مثبت، ۴ را مثبت و از ۵ نمونه کلاس منفی، ۴ را منفی تشخیص داده است. با آزمون دقیق فیشر یک‌طرفه بررسی کنید که آیا عملکرد آن به طور معنادار بهتر از حدس تصادفی است.
\end{exercise}
\begin{solution}
جدول: سطر «واقعاً مثبت» $(4,1)$ و سطر «واقعاً منفی» $(1,4)$؛ همه جمع‌های حاشیه‌ای $5$ و $n=10$. تحت $H_0$ (پیش‌بینی مستقل از برچسب واقعی)، تعداد پیش‌بینی‌های مثبت درست $a$ فوق‌هندسی است:
\[
\Prob(a=4)=\frac{\binom54\binom51}{\binom{10}5}=\frac{25}{252}=0.099,\qquad\Prob(a=5)=\frac{1}{252}=0.004 .
\]
$\text{p}=\Prob(a\ge4)=0.103>0.05$. با وجود دقت ۸۰٪، با فقط ۱۰ نمونه شواهد کافی برای ادعای «بهتر از تصادف» نداریم. درس: دقت بالا روی داده آزمون کوچک می‌تواند صرفاً شانس باشد.
\end{solution}

\begin{exercise}[انتخاب آزمون]
برای هر موقعیت، آزمون مناسب را نام ببرید: (الف) میانگین زمان پاسخ یک سرور با مقدار ادعایی مقایسه شود، $\sigma$ نامعلوم و $n=15$. (ب) آیا توزیع گروه خونی در یک شهر با نسبت‌های کشوری یکسان است؟ (ج) آیا دو الگوریتم روی ۸ مجموعه داده یکسان عملکرد متفاوتی دارند؟ (د) آیا زمان اجرای دو نسخه نرم‌افزار (با توزیعی بسیار چوله و نامعلوم) متفاوت است؟
\end{exercise}
\begin{solution}
(الف) آزمون $t$ تک‌نمونه‌ای (با فرض تقریباً گاوسی بودن داده‌ها). (ب) آزمون $\chi^2$ بهنجارش با $k$ = تعداد گروه‌های خونی و $p_i$ = نسبت‌های کشوری. (ج) آزمون $t$ \textbf{جفتی} روی تفاضل عملکرد در هر مجموعه داده (داده‌ها جفت‌اند). (د) آزمون جایگشت روی تفاضل میانگین (یا میانه)، چون فرض گاوسی بودن با داده چوله و نمونه کوچک قابل اعتماد نیست.
\end{solution}

\begin{exercise}[آزمون نسبت درست‌نمایی برای نمایی]
طول عمر قطعات $\mathrm{Exp}(\lambda)$ است. با لم نیمن--پیرسون، بهترین آزمون برای $H_0:\ \lambda=1$ در برابر $H_1:\ \lambda=2$ را بیابید.
\end{exercise}
\begin{solution}
$\log\Lambda_n=\sum_i\big[\log(2e^{-2X_i})-\log(e^{-X_i})\big]=n\log2-\sum_iX_i$. این عبارت با $\sum_iX_i$ نزولی است؛ پس «$\Lambda_n>\eta$» معادل «$\sum_iX_i<c$» است: اگر مجموع (یا میانگین) طول عمرها \textbf{کوچک} باشد $H_0$ رد می‌شود. منطقی است: $\lambda$ بزرگ‌تر یعنی میانگین عمر $1/\lambda$ کوتاه‌تر. (آستانه $c$ از توزیع $\sum_iX_i$ تحت $H_0$ که گاما است، یا برای $n$ بزرگ با حد مرکزی، تعیین می‌شود.)
\end{solution}
